% Definições

\newtheorem{def:base-modelos}{Definição}[chapter]
\newtheorem{def:uniforme discreta}[def:base-modelos]{Definição}
\newtheorem{def:modelo binomial}[def:base-modelos]{Definição}
\newtheorem{def:bernoulli}[def:base-modelos]{Definição}
\newtheorem{def:hipergeometrica}[def:base-modelos]{Definição}
\newtheorem{def:geometrica}[def:base-modelos]{Definição}
\newtheorem{def:binomial negativa}[def:base-modelos]{Definição}
\newtheorem{def:poisson}[def:base-modelos]{Definição}
\newtheorem{def:zeta}[def:base-modelos]{Definição}
\newtheorem{def:normal}[def:base-modelos]{Definição}

% Proposições
\newtheorem{prop:base-modelos}{Proposição}[chapter]
\newtheorem{prop:esperanca bernoulli}[prop:base-modelos]{Proposição}
\newtheorem{prop:esperanca binomial}[prop:base-modelos]{Proposição}
\newtheorem{prop:esperanca hipergeometrica}[prop:base-modelos]{Proposição}
\newtheorem{prop:esperanca geometrica}[prop:base-modelos]{Proposição}
\newtheorem{prop:esperanca binomial negativa}[prop:base-modelos]{Proposição}
\newtheorem{prop:esperanca poisson}[prop:base-modelos]{Proposição}
\newtheorem{prop:densidade da normal}[prop:base-modelos]{Proposição}
\newtheorem{prop:simetria densidade normal}[prop:base-modelos]{Proposição}
\newtheorem{prop:variancia 0 e igual media}[prop:base-modelos]{Proposição}

% Lema
\newtheorem{lem:base-modelos}{Lema}[chapter]
\newtheorem{lem:integral exponencial quadrada}[lem:base-modelos]{Lema}

% Teorema
\newtheorem{thm:base-modelos}{Teorema}[chapter]
\newtheorem{thm:tranformacao afim da normal}[thm:base-modelos]{Teorema}
\newtheorem{thm:parametros da normal}[thm:base-modelos]{Teorema}
\newtheorem{thm:simetria distribuicao normal}[thm:base-modelos]{Teorema}
\newtheorem{thm:demoivre-laplace}[thm:base-modelos]{Teorema}
\newtheorem{thm:desigualdade de markov}[thm:base-modelos]{Teorema}
\newtheorem{thm:lei fraca dos grandes numeros}[thm:base-modelos]{Teorema}

% Corolário
\newtheorem{cor:base-modelos}{Corolário}[chapter]
\newtheorem{cor:normal e normal padrao}[cor:base-modelos]{Corolário}
\newtheorem{cor:desigualdade de chebyshev}[cor:base-modelos]{Corolário}

% Exemplos
\theoremstyle{definition}
\newtheorem{exp:base-modelos}{Exemplo}[chapter]
\newtheorem{exp:uniforme discreta 1}[exp:base-modelos]{Exemplo}
\newtheorem{exp:uniforme discreta 2}[exp:base-modelos]{Exemplo}
\newtheorem{exp:bernoulli 1}[exp:base-modelos]{Exemplo}
\newtheorem{exp:bernoulli 2}[exp:base-modelos]{Exemplo}
\newtheorem{exp:binomial 1}[exp:base-modelos]{Exemplo}
\newtheorem{exp:hipergeometrica 1}[exp:base-modelos]{Exemplo}
\newtheorem{exp:geometrica 1}[exp:base-modelos]{Exemplo}
\newtheorem{exp:esperanca geometrica 1}[exp:base-modelos]{Exemplo}

\newtheorem{obs:base-modelos}{Observação}[chapter]
\newtheorem{obs:hipergeometrica 1}[obs:base-modelos]{Observação}
\newtheorem{obs:hipergeometrica 2}[obs:base-modelos]{Observação}
